\documentclass[border=5pt]{standalone}
\usepackage{tikz}
\usetikzlibrary{arrows.meta, calc, decorations.markings}

\begin{document}

\definecolor{axisX}{RGB}{220, 50, 30}   % Red - X forward (North)
\definecolor{axisY}{RGB}{40, 160, 40}   % Green - Y right (East)
\definecolor{axisZ}{RGB}{0, 120, 200}   % Blue - Z down
\definecolor{bodycolor}{RGB}{80, 80, 80}

\begin{tikzpicture}[
    axis/.style={-{Stealth[length=5mm, width=3mm]}, line width=2.5pt},
    rotarrow/.style={
      -{Stealth[length=3.5mm, width=2.5mm]},
      line width=1.8pt,
    },
  ]

  % ============================================================
  % True orthographic (parallel) projection of the FRD body frame
  % X=forward, Y=right, Z=down (right-handed).
  %
  % Camera model: view the body from behind (-X side), from the
  % right (+Y side), and from above (-Z side, since "up" in the
  % real world is the -Z direction here). This is parametrised by
  % an elevation angle ElevDeg (how far above the horizontal plane
  % the camera sits) and an azimuth AzimDeg (how far around from
  % "directly behind" toward "directly to the right" the camera
  % sits).
  %
  % Camera view direction (unit vector, camera -> scene):
  %   f = ( cos(el) cos(az), -cos(el) sin(az), sin(el) )
  % World "up" reference (real-world up, since Z points down):
  %   up_ref = (0, 0, -1)
  % Screen-right and screen-up unit vectors (standard look-at
  % construction, right = f x up_ref, up = right x f):
  %   r_hat = ( sin(az), cos(az), 0 )
  %   u_hat = ( cos(az) sin(el), -sin(az) sin(el), -cos(el) )
  % These closed forms were derived and their orthonormality
  % (r_hat . u_hat = 0, both unit length) verified numerically in
  % Python (P P^T = I to 1e-17) before being hard-coded below as
  % sin/cos of the two named angles -- no unexplained numbers.
  %
  % 正しい平行投影（直交投影）で FRD 機体座標系
  % （X=前, Y=右, Z=下, 右手系）を描く。
  %
  % カメラは機体の「後ろ（-X側）・右（+Y側）・上（この座標系では
  % 実世界の「上」は -Z 方向）」から見下ろす配置とし、仰角 ElevDeg
  % （水平面からどれだけ上に見上げる/見下ろすか）と方位角 AzimDeg
  % （「真後ろ」から「真右」に向けてどれだけ回り込むか）の2角で
  % 表す。
  %
  % カメラ視線方向（カメラ→対象、単位ベクトル）:
  %   f = ( cos(el) cos(az), -cos(el) sin(az), sin(el) )
  % 実世界の「上」基準ベクトル（Zが下なので (0,0,-1)）:
  %   up_ref = (0, 0, -1)
  % 画面右方向・画面上方向の単位ベクトル（標準的なlook-at構成:
  % right = f × up_ref, up = right × f）:
  %   r_hat = ( sin(az), cos(az), 0 )
  %   u_hat = ( cos(az) sin(el), -sin(az) sin(el), -cos(el) )
  % これらの閉形式はPythonで導出し、正規直交性
  % （r_hat・u_hat = 0、共に単位長）を P P^T = I (誤差1e-17) として
  % 数値検証済み。以下ではその2角のsin/cosとして直接埋め込む
  % （由来不明の数値は使わない）。
  \pgfmathsetmacro{\ElevDeg}{30}   % camera elevation above horizontal / 仰角
  \pgfmathsetmacro{\AzimDeg}{35}   % camera azimuth from "directly behind" / 方位角

  % Screen-projected components of the three world unit vectors,
  % i.e. the two rows of the 2x3 orthonormal projection matrix
  % evaluated on X=(1,0,0), Y=(0,1,0), Z=(0,0,1).
  % 世界座標系の3つの単位ベクトルを投影した画面上の成分
  % （2×3正規直交投影行列を X=(1,0,0), Y=(0,1,0), Z=(0,0,1) に
  % 適用した結果）。
  \pgfmathsetmacro{\Xx}{sin(\AzimDeg)}
  \pgfmathsetmacro{\Xy}{cos(\AzimDeg)*sin(\ElevDeg)}
  \pgfmathsetmacro{\Yx}{cos(\AzimDeg)}
  \pgfmathsetmacro{\Yy}{-sin(\AzimDeg)*sin(\ElevDeg)}
  \pgfmathsetmacro{\Zx}{0}
  \pgfmathsetmacro{\Zy}{-cos(\ElevDeg)}

  % Because an orthographic projection is a LINEAR map, projecting
  % any point built from linear combinations of X, Y, Z (body
  % corners, rotor circles, rotation-arc parametrisations) can be
  % done by taking the same linear combination of the already
  % projected \Xx/\Xy, \Yx/\Yy, \Zx/\Zy vectors below -- there is
  % no need to re-project every point from scratch.
  % 直交投影は線形写像なので、X, Y, Z の線形結合で作られる点
  % （機体コーナー、ローター円、回転円弧のパラメータ表示）は、
  % 個別に再投影せず、上で求めた投影済みベクトル \Xx/\Xy 等を
  % 同じ線形結合で使えばよい。

  % Geometry constants (all named, no bare magic numbers)
  % 幾何パラメータ（マジックナンバー無し、全て名前付き）
  \pgfmathsetmacro{\AL}{4.5}          % axis arrow length / 軸矢印の長さ
  \pgfmathsetmacro{\BodyHalf}{1.15}   % body plate half-diagonal / 機体プレート半対角長
  \pgfmathsetmacro{\RotorR}{0.5}      % rotor disk radius / ローター円半径
  \pgfmathsetmacro{\NoseLen}{0.4}     % front marker length beyond edge / 前方マーカーの突出量
  \pgfmathsetmacro{\NoseHalfW}{0.4}   % front marker half width / 前方マーカーの半幅
  % Rotation-arc placement was chosen, and verified numerically in
  % Python, so that the 270-degree circle (radius ArcR, centered at
  % ArcCenterDist along the axis) stays clear of all four rotor
  % disks (closest approach margin > 0.15 world units at these
  % values) -- an earlier smaller ArcCenterDist made the Roll arc
  % overlap the front-right rotor disk.
  % 回転円弧の配置は、270°の円（半径ArcR、軸上ArcCenterDistの位置が
  % 中心）が4つのローター円盤すべてから離れて見えるようPythonで
  % 数値検証して決定（この値で最近接マージン0.15ワールド単位超）。
  % 以前の小さいArcCenterDistではRoll弧が右前方ローターと重なっていた。
  \pgfmathsetmacro{\ArcCenterDist}{3.2} % rotation-arc center distance along axis / 回転円弧中心の軸上距離
  \pgfmathsetmacro{\ArcR}{0.65}        % rotation-arc radius / 回転円弧の半径

  % ---- Body plate: square in the XY (horizontal) plane, corners
  % at world coordinates (+-BodyHalf, +-BodyHalf, 0). FR/FL/RL/RR
  % name the four corners (Front/Rear, Right/Left).
  % ---- 機体プレート: XY平面（水平面）上の正方形。コーナーは世界
  % 座標で (±BodyHalf, ±BodyHalf, 0)。FR/FL/RL/RRは前後左右のコーナー名。
  \pgfmathsetmacro{\FRx}{\BodyHalf*\Xx+\BodyHalf*\Yx}
  \pgfmathsetmacro{\FRy}{\BodyHalf*\Xy+\BodyHalf*\Yy}
  \pgfmathsetmacro{\FLx}{\BodyHalf*\Xx-\BodyHalf*\Yx}
  \pgfmathsetmacro{\FLy}{\BodyHalf*\Xy-\BodyHalf*\Yy}
  \pgfmathsetmacro{\RLx}{-\BodyHalf*\Xx-\BodyHalf*\Yx}
  \pgfmathsetmacro{\RLy}{-\BodyHalf*\Xy-\BodyHalf*\Yy}
  \pgfmathsetmacro{\RRx}{-\BodyHalf*\Xx+\BodyHalf*\Yx}
  \pgfmathsetmacro{\RRy}{-\BodyHalf*\Xy+\BodyHalf*\Yy}

  \fill[bodycolor!10] (\FRx,\FRy) -- (\FLx,\FLy) -- (\RLx,\RLy) -- (\RRx,\RRy) -- cycle;
  \draw[bodycolor, line width=1.2pt] (\FRx,\FRy) -- (\FLx,\FLy) -- (\RLx,\RLy) -- (\RRx,\RRy) -- cycle;

  % ---- Rotor disks: true circles of radius RotorR lying in the XY
  % plane, centered on each body corner. Parametrised in 3D as
  % corner + RotorR*(cos(t) X + sin(t) Y) and projected by taking
  % the same combination of the projected \Xx/\Xy, \Yx/\Yy vectors
  % -- the projection is an ellipse, exactly as a circle on a
  % tilted plane appears under parallel projection.
  % ---- ローター円盤: XY平面上にある半径RotorRの真円。各機体コーナー
  % を中心とする。3D では corner + RotorR*(cos(t) X + sin(t) Y) と
  % パラメータ表示し、投影済みの \Xx/\Xy, \Yx/\Yy を同じ組み合わせで
  % 使って投影する（平行投影された傾いた平面上の円は楕円になる）。
  \foreach \Ccx/\Ccy in {\FRx/\FRy, \FLx/\FLy, \RLx/\RLy, \RRx/\RRy} {
    \draw[bodycolor, thin, fill=bodycolor!8]
      ({\Ccx+\RotorR*\Xx*cos(0)+\RotorR*\Yx*sin(0)},
       {\Ccy+\RotorR*\Xy*cos(0)+\RotorR*\Yy*sin(0)})
      \foreach \a in {6,12,...,360} {
        -- ({\Ccx+\RotorR*\Xx*cos(\a)+\RotorR*\Yx*sin(\a)},
            {\Ccy+\RotorR*\Xy*cos(\a)+\RotorR*\Yy*sin(\a)})
      } -- cycle;
  }

  % ---- Three body axes, drawn from the origin (body center) ----

  % X axis (Forward)
  \draw[axis, axisX] (0,0) -- ({\AL*\Xx},{\AL*\Xy});
  \node[font=\Large\bfseries, axisX, above right, inner sep=2pt]
    at ({\AL*\Xx},{\AL*\Xy}) {X (Forward)};

  % Y axis (Right)
  \draw[axis, axisY] (0,0) -- ({\AL*\Yx},{\AL*\Yy});
  \node[font=\Large\bfseries, axisY, below right, inner sep=2pt]
    at ({\AL*\Yx},{\AL*\Yy}) {Y (Right)};

  % Z axis (Down)
  \draw[axis, axisZ] (0,0) -- ({\AL*\Zx},{\AL*\Zy});
  \node[font=\Large\bfseries, axisZ, right, inner sep=2pt]
    at ({\AL*\Zx},{\AL*\Zy}) {Z (Down)};

  % ---- Front marker: small triangle on the front edge (FR-FL),
  % base at world (BodyHalf, +-NoseHalfW, 0), apex projecting
  % further forward at world (BodyHalf+NoseLen, 0, 0). Its base
  % midpoint sits exactly on the X-axis line (both are points along
  % world +X), so it is drawn AFTER the X axis: the marker's opaque
  % fill then paints over the line at that crossing, keeping the
  % marker readable as one solid shape instead of being bisected by
  % the axis stroke.
  % ---- 前方マーカー: 前縁(FR-FL)上の小さな三角形。底辺は世界座標で
  % (BodyHalf, ±NoseHalfW, 0)、頂点はさらに前方の
  % (BodyHalf+NoseLen, 0, 0) に投影される。底辺の中点はちょうど
  % X軸の直線上にある（どちらも世界+X上の点のため）ので、X軸を
  % 描いた後に描く。マーカーの不透明な塗りがその交点で線の上を
  % 上塗りし、軸線でマーカーが分断されずに1つの塗りつぶし形状として
  % 読めるようにする。
  \fill[axisX!75]
    ({\BodyHalf*\Xx-\NoseHalfW*\Yx}, {\BodyHalf*\Xy-\NoseHalfW*\Yy}) --
    ({\BodyHalf*\Xx+\NoseHalfW*\Yx}, {\BodyHalf*\Xy+\NoseHalfW*\Yy}) --
    ({(\BodyHalf+\NoseLen)*\Xx}, {(\BodyHalf+\NoseLen)*\Xy}) -- cycle;

  % === Rotation arrows =========================================
  % Each is a true 3D circle of radius ArcR, centered at
  % ArcCenterDist along the rotation axis, lying in the plane
  % perpendicular to that axis (spanned by the other two axes).
  % Parametrisation follows the cyclic right-hand rule
  % X -> Y -> Z -> X: for axis A with remaining axes (B, C) in that
  % cyclic order, a positive (right-hand) rotation about A is drawn
  % as center + ArcR*(cos(t) B + sin(t) C) for increasing t.
  % Each arc spans 270 degrees, leaving a 90 degree gap, and is
  % projected by taking the same linear combination of the already
  % projected axis vectors (valid because projection is linear).
  %
  % 各回転矢印は、回転軸に垂直な平面（残り2軸が張る面）内にある、
  % 軸上の点（原点からArcCenterDist）を中心とする半径ArcRの真円。
  % 巡回則 X→Y→Z→X に従い、軸Aの正（右手則）回転は、残り2軸を
  % (B, C) として center + ArcR*(cos(t) B + sin(t) C) （tは増加方向）
  % で描く。各弧は270°（90°の隙間を残す）。投影済みの軸ベクトルを
  % 同じ線形結合で使って投影する（投影が線形写像であるため妥当）。

  % Label placement rule (applies to all three arcs below): each
  % label is anchored at a point that is (ArcR + LabelGap) from the
  % arc center, in the radial direction of the arc's own MIDPOINT
  % angle (halfway between Start and End). This keeps every label
  % outside the stroke by a fixed clearance and roughly centered
  % over the visible (drawn) part of the arc, regardless of which
  % way that arc's ellipse happens to be oriented on screen.
  % ラベル配置規則（以下の3つの弧すべてに適用）: 各ラベルは、弧の
  % 中心から見て弧の中間角度（StartとEndの中間）の動径方向に
  % (ArcR + LabelGap) だけ離れた点に配置する。弧の楕円が画面上で
  % どの向きになっていても、線から一定の余白を保ちつつ、見えている
  % 弧のほぼ真上にラベルが来るようにするため。
  \pgfmathsetmacro{\LabelGap}{0.45}

  % ---- Roll: rotation about X axis, arc in the Y-Z plane (B=Y, C=Z)
  % ---- ロール: X軸まわりの回転、弧はYZ平面内 (B=Y, C=Z)
  \pgfmathsetmacro{\RollCx}{\ArcCenterDist*\Xx}
  \pgfmathsetmacro{\RollCy}{\ArcCenterDist*\Xy}
  \pgfmathsetmacro{\RollStart}{200}
  \pgfmathsetmacro{\RollEnd}{470}
  \pgfmathsetmacro{\RollMid}{(\RollStart+\RollEnd)/2}
  \draw[rotarrow, axisX!70!black]
    ({\RollCx+\ArcR*\Yx*cos(\RollStart)+\ArcR*\Zx*sin(\RollStart)},
     {\RollCy+\ArcR*\Yy*cos(\RollStart)+\ArcR*\Zy*sin(\RollStart)})
    \foreach \a in {206,212,...,470} {
      -- ({\RollCx+\ArcR*\Yx*cos(\a)+\ArcR*\Zx*sin(\a)},
          {\RollCy+\ArcR*\Yy*cos(\a)+\ArcR*\Zy*sin(\a)})
    };
  \node[font=\normalsize\bfseries, axisX!70!black, inner sep=2pt]
    at ({\RollCx+(\ArcR+\LabelGap)*\Yx*cos(\RollMid)+(\ArcR+\LabelGap)*\Zx*sin(\RollMid)},
        {\RollCy+(\ArcR+\LabelGap)*\Yy*cos(\RollMid)+(\ArcR+\LabelGap)*\Zy*sin(\RollMid)}) {Roll};

  % ---- Pitch: rotation about Y axis, arc in the Z-X plane (B=Z, C=X)
  % ---- ピッチ: Y軸まわりの回転、弧はZX平面内 (B=Z, C=X)
  \pgfmathsetmacro{\PitchCx}{\ArcCenterDist*\Yx}
  \pgfmathsetmacro{\PitchCy}{\ArcCenterDist*\Yy}
  \pgfmathsetmacro{\PitchStart}{200}
  \pgfmathsetmacro{\PitchEnd}{470}
  \pgfmathsetmacro{\PitchMid}{(\PitchStart+\PitchEnd)/2}
  \draw[rotarrow, axisY!70!black]
    ({\PitchCx+\ArcR*\Zx*cos(\PitchStart)+\ArcR*\Xx*sin(\PitchStart)},
     {\PitchCy+\ArcR*\Zy*cos(\PitchStart)+\ArcR*\Xy*sin(\PitchStart)})
    \foreach \a in {206,212,...,470} {
      -- ({\PitchCx+\ArcR*\Zx*cos(\a)+\ArcR*\Xx*sin(\a)},
          {\PitchCy+\ArcR*\Zy*cos(\a)+\ArcR*\Xy*sin(\a)})
    };
  \node[font=\normalsize\bfseries, axisY!70!black, inner sep=2pt]
    at ({\PitchCx+(\ArcR+\LabelGap)*\Zx*cos(\PitchMid)+(\ArcR+\LabelGap)*\Xx*sin(\PitchMid)},
        {\PitchCy+(\ArcR+\LabelGap)*\Zy*cos(\PitchMid)+(\ArcR+\LabelGap)*\Xy*sin(\PitchMid)}) {Pitch};

  % ---- Yaw: rotation about Z axis, arc in the X-Y plane (B=X, C=Y)
  % ---- ヨー: Z軸まわりの回転、弧はXY平面内 (B=X, C=Y)
  \pgfmathsetmacro{\YawCx}{\ArcCenterDist*\Zx}
  \pgfmathsetmacro{\YawCy}{\ArcCenterDist*\Zy}
  \pgfmathsetmacro{\YawStart}{160}
  \pgfmathsetmacro{\YawEnd}{430}
  \pgfmathsetmacro{\YawMid}{(\YawStart+\YawEnd)/2}
  \draw[rotarrow, axisZ!70!black]
    ({\YawCx+\ArcR*\Xx*cos(\YawStart)+\ArcR*\Yx*sin(\YawStart)},
     {\YawCy+\ArcR*\Xy*cos(\YawStart)+\ArcR*\Yy*sin(\YawStart)})
    \foreach \a in {166,172,...,430} {
      -- ({\YawCx+\ArcR*\Xx*cos(\a)+\ArcR*\Yx*sin(\a)},
          {\YawCy+\ArcR*\Xy*cos(\a)+\ArcR*\Yy*sin(\a)})
    };
  \node[font=\normalsize\bfseries, axisZ!70!black, inner sep=2pt]
    at ({\YawCx+(\ArcR+\LabelGap)*\Xx*cos(\YawMid)+(\ArcR+\LabelGap)*\Yx*sin(\YawMid)},
        {\YawCy+(\ArcR+\LabelGap)*\Xy*cos(\YawMid)+(\ArcR+\LabelGap)*\Yy*sin(\YawMid)}) {Yaw};

\end{tikzpicture}
\end{document}
